\comoleerla{(a) error porcentual mediano del pronóstico recursivo lanzado al inicio de la prueba, por horizonte, frente a las dos referencias; (b) aumento del error al permutar cada variable.}
{el error del modelo pasa de 4.5\,\% a un mes a 6.2\,\% a doce; es el menor de los tres en 7 de los 12 horizontes. La variable que más pesa es el tránsito del mes anterior.}
{al encadenar predicciones el error se acumula, y a varios meses la referencia estacional alcanza al modelo; por eso el pronóstico de la red se rotula como tal.}
